% !TEX TS-program = lualatex
% !TEX encoding = UTF-8

\documentclass[antiphonaire_ordinaire.tex]{subfiles}

\ifcsname preamble@file\endcsname
  \setcounter{page}{\getpagerefnumber{M-ao_communia}}
\fi

\begin{document}

\chapter*{Aux vigiles des apôtres}

\rubrica{Oraison propre. À défaut, ou si l'oraison \normaltext{Da quǽsumus} est assignée à la vigile mais déjà employée le même jour dans l'office ou la commémoration d'un confesseur pontife, on emploie l'oraison suivante.}

\oratio{Quǽsumus, omnípotens Deus:\pscross{} ut beátus \rubrica{N.} Apóstolus, cujus prævenímus festivitátem, tuum pro nobis impóret auxílium;\psstar{} ut a nostris reátibus absolúti, a cunctis étiam perículis eruámur. \perdominumla}{Nous vous en prions, Dieu tout-puissant, faites que le bienheureux Apôtre \rubrica{N.} dont nous anticipons la fête, implore pour nous votre secours, afin qu’étant purifiés de nos fautes, nous soyons aussi arrachés à tous les périls. \perdominumfr}

%%%%%%%%%%%%%%%%%%%%%%%%%%%%%%%%%%%%%%%%%%%

\chapter*{Commun des apôtres}

\section*{À Laudes}

\rubrica{Capitule, Hymne et Verset des Vêpres, \pageref{APEXVH}.}

\gscore{APEXLAB}{ad Ben.}{}

\section*{À Tierce}

\capitulum{Eph 2:19-20}{Fratres: Jam non estis hóspites et ádvenæ:\pscross{} sed estis cives Sanctórum et doméstici Dei: superædificáti super fundaméntum Apostolórum, et Prophetárum,\psstar{} ipso summo angulári lápide Christo Iesu.}{Frères, vous n’êtes plus des étrangers ni des gens de passage, vous êtes concitoyens des saints, vous êtes membres de la famille de Dieu, car vous avez été intégrés dans la construction qui a pour fondations les Apôtres et les prophètes ; et la pierre angulaire, c’est le Christ Jésus lui-même.}

\gscore{APEXTR}{}{\rr Leur bruit s’est répandu dans toute la terre. \vv Et leurs paroles jusqu’aux confins du globe de la terre.}

\versiculus{Constítues eos príncipes super omnem terram.}{Mémores erunt nóminis tui, Dómine.}{Vous les établirez princes sur toute la terre.}{Ils se souviendront de votre nom, Seigneur.}

\section*{À Sexte}

\capitulum{Ac. 5: 12}{Per manus autem Apostolórum\psstar{} fiébant signa et prodígia multa in plebe.}{Par les mains des Apôtres, beaucoup de signes et de prodiges s’accomplissaient dans le peuple.}

\gscore{APEXSR}{}{\rr Vous les établirez princes sur toute la terre. \vv Ils se souviendront de votre nom, Seigneur.}

\versiculus{Nimis honoráti sunt amíci tui Deus.}{Nimis confortátus est principátus eórum.}{Vos amis ont été grandement honorés, ô Dieu.}{Leur autorité de princes a été établie avec puissance.}

\section*{À None}

\capitulum{Ac. 5: 41}{Ibant Apóstoli gaudéntes a conspéctu concílii,\psstar{} quóniam digni hábiti sunt pro nómine Jesu contuméliam pati.}{Les Apôtres, quittant le Conseil suprême, repartaient tout joyeux d’avoir été jugés dignes de subir des humiliations pour le nom de Jésus.}

\gscore{APEXNR}{}{\rr Vos amis ont été grandement honorés, ô Dieu. \vv Leur autorité de princes a été établie avec puissance.}

\versiculus{Annuntiavérunt ópera Dei.}{Et facta ejus intellexérunt.}{Ils ont annoncé les œuvres de Dieu.}{Et ils ont compris les choses qu’il a faites.}

\section*{À Vêpres}

\capitulum{Eph 2:19-20}{Fratres: Jam non estis hóspites et ádvenæ:\pscross{} sed estis cives Sanctórum et doméstici Dei: superædificáti super fundaméntum Apostolórum, et Prophetárum,\psstar{} ipso summo angulári lápide Christo Iesu.}{Frères, vous n’êtes plus des étrangers ni des gens de passage, vous êtes concitoyens des saints, vous êtes membres de la famille de Dieu, car vous avez été intégrés dans la construction qui a pour fondations les Apôtres et les prophètes ; et la pierre angulaire, c’est le Christ Jésus lui-même.}

\hymnus{APEXVH}
\gscore{APEXVHb}{}{}
% TODO Hymni antiqui

\versiculus{Annuntiavérunt ópera Dei.}{Et facta ejus intellexérunt.}{Ils ont annoncé les œuvres de Dieu.}{Et ils ont compris les choses qu’il a faites.}

\gscore{APEX2VAM}{ad Mag.}{Soyez courageux dans la guerre ; combattez contre l’ancien serpent, et vous recevrez le royaume éternel. Allelúia.}

\rubrica{Oraison propre.}

%%%%%%%%%%%%%%%%%%%%%%%%%%%%%%%%%%%%%%%%%%%

\chapter*{Commun d'un martyr}

\section*{À Laudes}

\capitulum{Jc. 1: 12}{Beátus vir, qui suffert tentatiónem:\pscross{} quóniam cum probátus fúerit, accípiet corónam vitæ,\psstar{} quam repromísit Deus diligéntibus se.}{Heureux l’homme qui supporte l’épreuve avec persévérance, car, sa valeur une fois vérifiée, il recevra la couronne de la vie promise à ceux qui aiment Dieu.}

\hymnus{UMEXLHa}
\hymnus{UMEXLHb}

%% todo hymne antique

\versiculus{Justus ut palma florébit.}{Sicut cedrus Líbani multiplicábitur.}{Le juste fleurira comme le palmier.}{Il croîtra comme le cèdre du Liban.}

\gscore{UMEXLAB}{Ad Ben.}{Celui qui méprise son âme en ce monde, la conserve pour la vie éternelle.}

\rubrica{Oraison des Vêpres.}

\section*{À Tierce}

\capitulum{Jc. 1: 12}{Beátus vir, qui suffert tentatiónem:\pscross{} quóniam cum probátus fúerit, accípiet corónam vitæ,\psstar{} quam repromísit Deus diligéntibus se.}{Heureux l’homme qui supporte l’épreuve avec persévérance, car, sa valeur une fois vérifiée, il recevra la couronne de la vie promise à ceux qui aiment Dieu.}

\gscore{UMEXTR}{}{\rr De gloire et d'honneur, Seigneur, Vous l'avez couronné. \vv Et Vous lui avez donné l'empire sur les œuvres de Vos mains.}

\versiculus{Posuísti, Dómine, super caput ejus.}{Corónam de lápide pretióso.}{Vous avez mis, Seigneur, sur sa tête.}{Une couronne de pierres précieuses.}

\section*{À Sexte}

\capitulum{Si. 15: 3}{Cibávit illum pane vitæ et intelléctus,\pscross{} et aqua sapiéntiæ salutáris\psstar{} potávit illum Dóminus Deus noster.}{Le Seigneur notre Dieu le nourrit du pain de l’intelligence, et lui donne à boire l’eau de la sagesse.}

\gscore{UMEXSR}{}{\rr Vous avez mis Seigneur, sur sa tête. \vv Une couronne de pierres précieuses.}

\versiculus{Magna est glória ejus in salutári tuo.}{Glóriam et magnum decórem impónes super eum.}{Sa gloire est grande, grâce à Votre secours.}{Vous l'avez revêtu de splendeur et de magnificence.}

\section*{À None}

\capitulum{Si. 39: 5}{Justus cor suum trádidit ad vigilándum dilúculo ad Dóminum, qui fecit illum,\psstar{} et in conspéctu Altíssimi deprecábitur.}{Le juste s’appliquera de tout son cœur à servir dès le matin le Seigneur qui l’a créé, il présentera sa supplication devant le Très-Haut.}

\gscore{UMEXNR}{}{\rr Sa gloire est grande, grâce à Votre secours. \vv Vous l'avez revêtu de splendeur et de magnificence.}

\versiculus{Justus ut palma florébit.}{Sicut cedrus Líbani multiplicábitur.}{Le juste fleurira comme le palmier.}{Il croîtra comme le cèdre du Liban.}

\section*{À Vêpres}

\capitulum{Jc. 1: 12}{Beátus vir, qui suffert tentatiónem:\pscross{} quóniam cum probátus fúerit, accípiet corónam vitæ,\psstar{} quam repromísit Deus diligéntibus se.}{Heureux l’homme qui supporte l’épreuve avec persévérance, car, sa valeur une fois vérifiée, il recevra la couronne de la vie promise à ceux qui aiment Dieu.}

\hymnus{UMEXVH}
\hymnus{UMEXVHb}
% todo hymni antiqui

\versiculus{Justus ut palma florébit.}{Sicut cedrus Líbani multiplicábitur.}{Le juste fleurira comme le palmier.}{Il croîtra comme le cèdre du Liban.}

\gscore{UMEXVAM}{Ad Mag.}{Si quelqu'un veut marcher derrière moi, qu'il se renonce lui-même, qu'il porte sa croix et me suive.}

\rubrica{Pour un souverain pontife martyr:}

\oratio{Gregem tuum, Pastor ætérne, placátus inténde:\pscross{} et, per beátum \rubrica{N.} Mártyrem tuum atque Summum Pontíficem, perpétua protectióne custódi;\psstar{} quem totíus Ecclésiæ præstitísti esse pastórem. \perdominumla}{Pasteur éternel, considérez avec bonté votre troupeau, et gardez-le sous la constante protection du bienheureux \rubrica{N.}, votre martyr et souverain pontife, que vous avez constitué pasteur de toute l’Église. \perdominumfr}

\rubrica{Pour la commémoraison d'un souverain pontife, si l'oraison \normaltext{Gregem tuum} a déjà été employée:}

\oratio{Deus, qui Ecclésiam tuam, in apostólicæ petræ soliditáte fundátam, ab infernárum éruis terróre portárum:\pscross{} præsta, quǽsumus; ut, intercedénte beáto \rubrica{N.} Mártyre tuo atque Summo Pontífice, in tua veritáte persístens,\psstar{} contínua securitáte muniátur. \perdominumla}{Seigneur Dieu, ayant fondé votre Église sur la solide pierre apostolique, vous la sauvez de la terreur des puissances de l’enfer. Par l’intercession du saint pape \rubrica{N.} votre martyr, daignez accorder à votre Église de demeurer ferme dans votre vérité, assurée de votre continuelle protection. \perdominumfr}

\rubrica{Pour un martyr pontife:}

\oratio{Infirmitátem nostram réspice, omnípotens Deus:\pscross{} et, quia pondus própriæ actiónis gravat,\psstar{} beáti \rubr{N.} Mártyris tui atque Pontíficis intercéssio gloriósa nos prótegat. \perdominumla}{Dieu tout-puissant, regardez notre fai­blesse ; et parce que le poids de nos propres péchés nous accable, que la glorieuse intercession du bienheu­reux \rubr{N.}, votre évêque et martyr, nous protège. \perdominumfr}

\rubrica{Ou bien:}

\oratio{Deus, qui nos beáti \rubr{N.} Mártyris tui atque Pontíficis ánnua solemnitáte lætíficas:\pscross{} concéde propítius; ut, cujus natalítia cólimus,\psstar{} de ejúsdem étiam protectióne gaudeámus. \perdominumla}{Ô Dieu qui nous donnez la joie de la solennité annuelle de Votre bienheureux Martyr et Pontife \rubr{N.}, faites, dans Votre bonté, que célébrant sa naissance, nous nous réjouissions aussi de sa protection. \perdominumfr}

\rubrica{Pour un martyr non-pontife:}

\oratio{Præsta, quǽsumus, omnípotens Deus:\pscross{} ut, qui beáti \rubr{N.} Mártyris tui natalítia cólimus,\psstar{} intercessióne ejus in tui nóminis amóre roborémur. \perdominumla}{Faites, s’il Vous plaît, Dieu tout-puissant, que célébrant la naissance de Votre bienheureux Martyr \rubr{N.}, nous soyons, par son intercession fortifiés dans l’amour de Votre nom. \perdominumfr}

\rubrica{Ou bien:}

\oratio{Præsta, quǽsumus, omnípotens Deus:\pscross{} ut, intercedénte beáto \rubr{N.} Mártyre tuo, et a cunctis adversitátibus liberémur in córpore,\psstar{} et a pravis cogitatiónibus mundémur in mente. \perdominumla}{Faites, s’il Vous plaît, Dieu tout-puissant, que célébrant la naissance de Votre bienheureux Martyr \rubr{N.}, nos corps soient affranchis de toute adversité et nos âmes purifiées des pensées mauvaises. \perdominumfr}


%%%%%%%%%%%%%%%%%%%%%%%%%%%%%%%%%%%%%%%%%%%

\chapter*{Commun de plusieurs martyrs}

\section*{À Laudes}

\capitulum{Sg. 3: 1-3}{Justórum ánimæ in manu Dei sunt, et non tanget illos torméntum mortis.\pscross{} Visi sunt óculis insipiéntium mori:\psstar{} illi autem sunt in pace.}{Les âmes des justes sont dans la main de Dieu ; aucun tourment n’a de prise sur eux. Aux yeux de l’insensé, ils ont paru mourir : mais ils sont dans la paix.}

\hymnus{PMEXLH}
\hymnus{PMEXLHb}
% todo hymni antiqui

\versiculus{Exsultábunt Sancti in glória.}{Lætabúntur in cubílibus suis.}{Les Saints exulteront en gloire.}{Ils se réjouiront sur leurs lits de repos.}

\gscore{PMEXLAB}{Ad Ben.}{Les cheveux de votre tête sont tous comptés. Ne craignez pas ; vous valez mieux que beaucoup de passereaux.}

\rubrica{Oraison des Vêpres.}

\section*{À Tierce}

\capitulum{Sg. 3: 1-3}{Justórum ánimæ in manu Dei sunt, et non tanget illos torméntum mortis.\pscross{} Visi sunt óculis insipiéntium mori:\psstar{} illi autem sunt in pace.}{Les âmes des justes sont dans la main de Dieu ; aucun tourment n’a de prise sur eux. Aux yeux de l’insensé, ils ont paru mourir : mais ils sont dans la paix.}

\gscore{PMEXTR}{}{\rr Réjouissez-vous dans le Seigneur, et exultez, ô justes. \vv Et soyez glorifiés vous tous qui avez le coeur droit.}

\versiculus{Exsúltent justi in conspéctu Dei.}{Et delecténtur in lætítia.}{Qu’ils exultent, les justes, en présence de Dieu.}{Et qu’ils se délectent dans la joie.}

\section*{À Sexte}

\capitulum{Sg. 10: 17}{Réddidit Deus mercédem labórum sanctórum suórum,\pscross{} et dedúxit illos in via mirábili: et fuit illis in velaménto diéi,\psstar{} et in luce stellárum nocte.}{La Sagesse a récompensé les saints de leurs peines, les a conduits sur un chemin de merveilles. Le jour, elle fut pour eux un abri, et la nuit, une clarté d’étoiles.
}
\gscore{PMEXSR}{}{\rr Qu’ils exultent les justes, en présence de Dieu. \vv Et qu’ils se délectent dans la joie.}

\versiculus{Justi autem in perpétuum vivent.}{Et apud Dóminum est merces eórum.}{Quant aux justes, ils vivront éternellement.}{Et auprès du Seigneur est leur récompense.}

\section*{À None}

\capitulum{Sg. 3: 7-8}{Fulgébunt justi, et tamquam scintíllæ in arundinéto discúrrent.\pscross{} Judicábunt natiónes, et dominabúntur pópulis:\psstar{} et regnábit Dóminus illórum in perpétuum.}{Les justes resplendiront; comme l’étincelle qui court sur la paille, ils avancent. Ils jugeront les nations, ils auront pouvoir sur les peuples, et le Seigneur régnera sur eux pour les siècles.}

\gscore{PMEXNR}{}{\rr Quant aux justes, ils vivront éternellement. \vv Et auprès du Seigneur est leur récompense.}

\versiculus{Exsultábunt Sancti in glória.}{Lætabúntur in cubílibus suis.}{Les Saints exulteront en gloire.}{Ils se réjouiront sur leurs lits de repos.}

\section*{À Vêpres}

\capitulum{Sg. 3: 1-3}{Justórum ánimæ in manu Dei sunt, et non tanget illos torméntum mortis.\pscross{} Visi sunt óculis insipiéntium mori:\psstar{} illi autem sunt in pace.}{Les âmes des justes sont dans la main de Dieu ; aucun tourment n’a de prise sur eux. Aux yeux de l’insensé, ils ont paru mourir : mais ils sont dans la paix.}

\hymnus{PMEXVHa}
\hymnus{PMEXVHb}

% TODO hymni antiqui

\versiculus{Exsultábunt Sancti in glória.}{Lætabúntur in cubílibus suis.}{Les Saints exulteront en gloire.}{Ils se réjouiront sur leurs lits de repos.}

\gscore{PMEXVAM}{Ad Mag.}{Elles se réjouissent dans les cieux, les âmes des saints qui ont suivi les traces du Christ. Et, parce qu’ils ont répandu leur sang pour Son amour, avec le Christ, ils exultent sans fin.}


\rubrica{Pour plusieurs martyrs souverains pontifes}

\oratio{v. Gregem tuum, Pastor ætérne, placátus inténde:\pscross{} et per beátos \rubrica{N.} et \rubrica{N.} Mártyres tuos atque Summos Pontífices, perpétua protectióne custódi;\psstar{} quos totíus Ecclésiæ præstitísti esse pastóres. \perdominumla}{Pasteur éternel, considérez avec bonté votre troupeau, et gardez-le sous la constante protection des bienheureux \rubrica{N.} et \rubrica{N.}, vos martyrs et souverains pontifes, que vous avez constitués pasteurs de toute l’Église. \perdominumfr}

\rubrica{Pour plusieurs martyrs pontifes}

\oratio{Beatórum Mártyrum paritérque Pontíficum \rubr{N.} et \rubr{N.} nos, quǽsumus, Dómine, festa tueántur:\psstar{} et eórum comméndet orátio veneránda. \perdominumla}{Que les fêtes des bienheureux martyrs et pontifes \rubr{N.} et \rubr{N.} nous protègent, Seigneur, s'il Vous plaît ; et que leur précieuse intercession nous soutienne. \perdominumfr}

\rubrica{Pour plusieurs martyrs non-pontifes}

\oratio{Deus, qui nos concédis sanctórum Mártyrum tuórum \rubr{N.} et \rubr{N.} natalítia cólere:\pscross{} da nobis in ætérna beatitúdine\psstar{} de eórum societáte gaudére. \perdominumla}{Dieu, qui nous permettez de cé­lébrer la naissance au ciel de vos saints martyrs \rubr{N.} et \rubr{N.}, accordez-nous de jouir de leur société dans l’éternité bienheureuse. \perdominumfr}

\rubrica{Ou bien:}

\oratio{Deus, qui nos ánnua sanctórum Mártyrum tuórum \rubr{N.} et \rubr{N.} solemnitáte lætíficas:\pscross{} concéde propítius; ut quorum gaudémus méritis,\psstar{} accendámur exémplis. \perdominumla}{Ô Dieu, qui nous réjouissez par la solennité de Vos bienheureux \rubr{N.} et \rubr{N.}, accordez-nous, s’il Vous plaît, la grâce d’être enflammés par les exemples de ceux dont les mérites nous réjouissent. \perdominumfr}

%%%%%%%%%%%%%%%%%%%%%%%%%%%%%%%%%%%%%%%%%%%

\chapter*{Commun d'un confesseur pontife}

\section*{À Laudes}

\capitulum{Si. 44: 16-17}{Ecce sacérdos magnus,\pscross{} qui in diébus suis plácuit Deo, et invéntus est justus:\psstar{} et in témpore iracúndiæ factus est reconciliátio.}{Voici le grand-prêtre qui, aux jours de sa vie, fut agréable au Seigneur, et fut trouvé juste; au temps de la colère, il a été l’instrument de la réconciliation.}

\hymnus{COPOLHa}
\hymnus{COPOLHb} % TODO hymni antiqui
\hymnus{COPOLHc}

\versiculus{Justum dedúxit Dóminus per vias rectas.}{Et osténdit illi regnum Dei.}{Le Seigneur l’a conduit par des voies de droiture.}{Et il lui a montré le royaume de Dieu.}

\gscore{}{Ad Ben.}{C’est bien, serviteur bon et fidèle, parce que tu as été fidèle sur peu de choses, je t’établirai sur beaucoup, dit le Seigneur.}

\rubrica{Oraison des Vêpres.}

\section*{À Tierce}

\capitulum{Si. 44: 16-17}{Ecce sacérdos magnus,\pscross{} qui in diébus suis plácuit Deo, et invéntus est justus:\psstar{} et in témpore iracúndiæ factus est reconciliátio.}{Voici le grand-prêtre qui, aux jours de sa vie, fut agréable au Seigneur, et fut trouvé juste; au temps de la colère, il a été l’instrument de la réconciliation.}

\gscore{COPOTR}{}{\rr Le Seigneur l’a aimé, et l’a paré. \vv Il l’a revêtu de la robe de gloire.}

\versiculus{Elégit eum Dóminus sacerdótem sibi.}{Ad sacrificándum ei hóstiam laudis.}{Le Seigneur l’a choisi pour son prêtre.}{Pour lui sacrifier l’hostie de louange.}

\section*{À Sexte}

\capitulum{Si. 44: 19, 21}{Non est invéntus símilis illi,\pscross{} qui conserváret legem Excélsi:\psstar{} ídeo jurejurándo fecit illum Dóminus créscere in plebem suam.}{Personne n’a jamais égalé sa gloire, car il observa la loi du Très-Haut; aussi Dieu lui a-t-il assuré par serment qu’il le multiplierait au milieu de son peuple.}

\gscore{COPOSR}{}{\rr Le Seigneur l’a choisi, pour être son prêtre. \vv Pour lui sacrifier l’hostie de louange.}

\versiculus{Tu es sacérdos in ætérnum.}{Secúndum órdinem Melchísedech.}{Tu es Prêtre pour l’éternité.}{Selon l’ordre de Melchisédech.}

\section*{À None}

\capitulum{Si. 45: 15-16}{Fungi sacerdótio,\pscross{} et habére laudem in nómine ipsíus,\psstar{} et offérre illi incénsum dignum in odórem suavitátis.}{Pour exercer le sacerdoce, bénir le peuple par le nom du Seigneur, et lui présenter l’offrande, l’encens et le parfum en mémorial.}

\gscore{COPONR}{}{\rr Tu es prêtre pour l’éternité. \vv Selon l’ordre de Melchisédech.}

\versiculus{Justum dedúxit Dóminus per vias rectas.}{Et osténdit illi regnum Dei.}{Le Seigneur l’a conduit par des voies de droiture.}{Et il lui a montré le royaume de Dieu.}

\section*{À Vêpres}

\capitulum{Si. 44: 16-17}{Ecce sacérdos magnus,\pscross{} qui in diébus suis plácuit Deo, et invéntus est justus:\psstar{} et in témpore iracúndiæ factus est reconciliátio.}{Voici le grand-prêtre qui, aux jours de sa vie, fut agréable au Seigneur, et fut trouvé juste; au temps de la colère, il a été l’instrument de la réconciliation.}

\hymnus{COPOVH1}
\hymnus{COPOVH2}
\hymnus{COPOVH3}
\hymnus{COPOVH4}

%todo hymni antiqui

\versiculus{Justum dedúxit Dóminus per vias rectas.}{Et osténdit illi regnum Dei.}{Le Seigneur l’a conduit par des voies de droiture.}{Et il lui a montré le royaume de Dieu.}

\rubrica{Pour un souverain pontife:}
\gscore{COSPVAM}{Ad Mag.}{Tandis qu'il était Souverain Pontife, il n'a craint rien de terrestre, mais il s'en est allé glorieux au royaume céleste.}

\rubrica{Pour un pontife:}
\gscore{COPOVAM}{Ad Mag.}{Le Seigneur l'a aimé, et l'a paré. Il l'a revêtu de la robe de gloire et l'a couronné aux portes du paradis.4
}

\rubrica{Pour un docteur pontife:}
\gscore{CODOVAM}{Ad Mag.}{O docteur excellent, lumière de la Sainte Eglise, bienheureux N., amoureux de la loi divine, priez pour nous le Fils de Dieu.}

\rubrica{Pour un confesseur souverain pontife:}

\oratio{Gregem tuum, Pastor ætérne, placátus inténde:\pscross{} et, per beátum \rubrica{N.} Summum Pontíficem, perpétua protectióne custódi;\psstar{} quem totíus Ecclésiæ præstitísti esse pastórem. \perdominumla}{Pasteur éternel, considérez avec bonté votre troupeau, et gardez-le sous la constante protection du bienheureux \rubrica{N.}, souverain pontife, que vous avez constitué pasteur de toute l’Église. \perdominumfr}

\rubrica{Pour la commémoraison d'un souverain pontife, si l'oraison \normaltext{Gregem tuum} a déjà été employée:}

\oratio{Deus, qui Ecclésiam tuam, in apostólicæ petræ soliditáte fundátam, ab infernárum éruis terróre portárum:\pscross{} præsta, quǽsumus; ut, intercedénte beáto \rubrica{N.} Summo Pontífice, in tua veritáte persístens,\psstar{} contínua securitáte muniátur. \perdominumla}{Seigneur Dieu, ayant fondé votre Église sur la solide pierre apostolique, vous la sauvez de la terreur des puissances de l’enfer. Par l’intercession du saint pape \rubrica{N.}, daignez accorder à votre Église de demeurer ferme dans votre vérité, assurée de votre continuelle protection. \perdominumfr}

\rubrica{Pour un confesseur pontife:}

\oratio{Da, quǽsumus, omnípotens Deus:\pscross{} ut beáti \rubr{N.} Confessóris tui atque Pontíficis veneránda solémnitas,\psstar{} et devotiónem nobis áugeat et salútem. \perdominumla}{Accordez à notre prière, Dieu tout-puissant, que la vénérable solennité de votre bienheureux \rubr{N.}, Confesseur et Pontife, augmente en nous dévotion et santé. \perdominumfr}

\rubrica{Ou bien:}

\oratio{Exáudi, quǽsumus, Dómine, preces nostras, quas in beáti \rubr{N.} Confessóris tui atque Pontíficis solemnitáte deférimus:\pscross{} et, qui tibi digne méruit famulári, ejus intercedéntibus méritis,\psstar{} ab ómnibus nos absólve peccátis. \perdominumla}{Exaucez, Seigneur, s’il Vous plaît, les prières qu’en la solennité du bienheureux \rubr{N.}, Votre Confesseur et Pontife, nous vous présentons et, par l’intercession des mérites de celui qui a obtenu de Vous servir dignement, absolvez-nous de tous nos péchés. \perdominumfr}

\rubrica{Pour un docteur pontife:}

\oratio{Deus, qui pópulo tuo ætérnæ salútis beátum \rubr{N.} minístrum tribuísti:\pscross{} præsta, quǽsumus; ut, quem Doctórem vitæ habúimus in terris,\psstar{} intercessórem habére mereámur in cœlis. \perdominumla}{Ô Dieu, qui avez donné à votre peuple le bienheureux \rubr{N.} comme ministre du salut éternel, nous vous en prions : faites que, l’ayant eu comme maître de vie sur la terre, nous méritions de l’avoir comme intercesseur dans le ciel. \perdominumfr}

\rubrica{Pour plusieurs confesseurs pontifes, là où une semblable fête existe, on dit \normaltext{beatórum N. et N. Confessórum tuum atque Pontíficum} à la place de \normaltext{beáti N. Confessóris tui atque Pontíficis}, et, dans le cas de souverains pontifes, \normaltext{beátos \rubrica{N.} et \rubrica{N.} Summos Pontífices} à la place de \normaltext{beátum \rubrica{N.} Summum Pontíficem}.}


%%%%%%%%%%%%%%%%%%%%%%%%%%%%%%%%%%%%%%%%%%%

\chapter*{Commun d'un confesseur non-pontife}

\section*{À Laudes}

\capitulum{Eccl. 31: 8-9}{Beátus vir, qui invéntus est sine mácula,\pscross{} et qui post aurum non ábiit, nec sperávit in pecúnia et thesáuris.\psstar{} Quis est hic, et laudábimus eum? fecit enim mirabília in vita sua.}{Heureux l'homme qui fut trouvé sans reproche, et n’a pas couru après l’or, et n'a pas mis son espoir dans l'argent et les trésors. Qui est-il ? Nous le dirons bienheureux : dans sa vie, il a fait des merveilles.}

\hymnus{CONPLHa}
\hymnus{CONPLHb}
\hymnus{CONPLHc} % TODO hymni antiqui

\versiculus{Justum dedúxit Dóminus per vias rectas.}{Et osténdit illi regnum Dei.}{Le Seigneur l’a conduit par des voies de droiture.}{Et il lui a montré le royaume de Dieu.}

\gscore{CONPLAB}{Ad Ben.}{Très bien, bon serviteur, fidèle en peu de choses, parce que tu as été fidèle en peu de choses, je t'établirai sur beaucoup, entre dans la joie de ton Seigneur.}

\rubrica{Oraison des Vêpres.}

\section*{À Tierce}

\capitulum{Eccl. 31: 8-9}{Beátus vir, qui invéntus est sine mácula,\pscross{} et qui post aurum non ábiit, nec sperávit in pecúnia et thesáuris.\psstar{} Quis est hic, et laudábimus eum? fecit enim mirabília in vita sua.}{Heureux l'homme qui fut trouvé sans reproche, et n’a pas couru après l’or, et n'a pas mis son espoir dans l'argent et les trésors. Qui est-il ? Nous le dirons bienheureux : dans sa vie, il a fait des merveilles.}

\gscore{COPOTR}{}{\rr Le Seigneur l’a aimé, et l’a paré. \vv Il l’a revêtu de la robe de gloire.} % repris de COPO, normal

\versiculus{Os justi meditábitur sapiéntiam.}{Et lingua ejus loquétur judícium.}{La bouche du juste redira la sagesse.}{Et sa langue publiera la justice.}

\section*{À Sexte}

\capitulum{Si. 39: 5}{Justus cor suum trádidit ad vigilándum dilúculo ad Dóminum, qui fecit illum,\psstar{} et in conspéctu Altíssimi deprecábitur.}{Le juste s’appliquera de tout son cœur à servir dès le matin le Seigneur qui l’a créé, et il présentera sa supplication devant le Très-Haut.}

\gscore{CONPSR}{}{\rr La bouche du juste redira la sagesse. \vv Et sa langue publiera la justice.}

\versiculus{Lex Dei ejus in corde ipsíus.}{Et non supplantabúntur gressus ejus.}{La loi de son Dieu est dans son cœur.}{Et il ne fera pas de faux pas.}

\section*{À None}

\capitulum{Sg. 10: 10}{Justum dedúxit Dóminus per vias rectas, et osténdit illi regnum Dei, et dedit illi sciéntiam sanctórum: honestávit illum in labóribus, et complévit labóres illíus.}{Le Seigneur guida le juste sur des sentiers droits : il lui fit voir le royaume de Dieu et lui donna la connaissance de réalités saintes ; il récompensa ses efforts et multiplia le fruit de ses labeurs.}

\gscore{CONPNR}{}{\rr La loi de son Dieu est dans son cœur. \vv Et il ne fera pas de faux pas.}

\versiculus{Justum dedúxit Dóminus per vias rectas.}{Et osténdit illi regnum Dei.}{Le Seigneur l’a conduit par des voies de droiture.}{Et il lui a montré le royaume de Dieu.}

\section*{À Vêpres}

\capitulum{Eccl. 31: 8-9}{Beátus vir, qui invéntus est sine mácula,\pscross{} et qui post aurum non ábiit, nec sperávit in pecúnia et thesáuris.\psstar{} Quis est hic, et laudábimus eum? fecit enim mirabília in vita sua.}{Heureux l'homme qui fut trouvé sans reproche, et n’a pas couru après l’or, et n'a pas mis son espoir dans l'argent et les trésors. Qui est-il ? Nous le dirons bienheureux : dans sa vie, il a fait des merveilles.}

\hymnus{CONPVH5}
\hymnus{CONPVH6}
\hymnus{CONPVH7}
\hymnus{CONPVH8}

\versiculus{Justum dedúxit Dóminus per vias rectas.}{Et osténdit illi regnum Dei.}{Le Seigneur l’a conduit par des voies de droiture.}{Et il lui a montré le royaume de Dieu.}

\gscore{CONPVAM}{Ad Mag.}{Cet homme, méprisant le monde et les choses de la terre, s’est assuré, triomphant, par sa parole et ses actes, des richesses dans le ciel.}

\oratio{Deus, qui nos beáti \rubr{N.} Confessóris tui ánnua solemnitáte lætíficas:\pscross{} concéde propítius; ut cujus natalítia cólimus,\psstar{} étiam actiónes imitémur. \perdominumla}{O Dieu, qui nous réjouissez par la fête annuelle du bienheureux \rubr{N.}, Votre confesseur, accordez-nous miséricordieusement que, fêtant sa naissance au ciel, nous imitions aussi ses actions. \perdominumfr}

\oratio{Adésto, Dómine, supplicatiónibus nostris, quas in beáti \rubr{N.} Confessóris tui sollemnitáte deférimus:\pscross{} ut, qui nostræ justítiæ fidúciam non habémus,\psstar{} ejus, qui tibi plácuit, précibus adjuvémur. \perdominumla}{Soyez favorable, Seigneur, aux supplications que nous vous présentons en la solennité du bienheureux \rubr{N.}, votre confesseur ; afin que, n’ayant pas confiance en notre justice, nous soyons aidés par les prières de celui qui vous a plu. \perdominumfr}

\rubrica{Pour un docteur}

\oratio{Deus, qui pópulo tuo ætérnæ salútis beátum \rubr{N.} minístrum tribuísti:\pscross{} præsta, quǽsumus; ut, quem Doctórem vitæ habúimus in terris,\psstar{} intercessórem habére mereámur in cœlis. \perdominumla}{Ô Dieu, qui avez donné à votre peuple le bienheureux \rubr{N.} comme ministre du salut éternel, nous vous en prions : faites que, l’ayant eu comme maître de vie sur la terre, nous méritions de l’avoir comme intercesseur dans le ciel. \perdominumfr}

\rubrica{Pour un abbé}

\oratio{Intercéssio nos, quǽsumus, Dómine, beáti \rubr{N.} Abbátis comméndet:\pscross{} ut, quod nostris méritis non valémus,\psstar{} ejus patrocínio assequámur. \perdominumla}{Que l’intercession du bienheureux abbé \rubr{N.} nous recommande auprès de vous, Seigneur, afin que son patronage nous obtienne ce que nous ne pouvons acquérir par nos mérites. \perdominumfr}

%%%%%%%%%%%%%%%%%%%%%%%%%%%%%%%%%%%%%%%%%%%

\chapter*{Commun des vierges}

\section*{À Laudes}

\capitulum{2 Co. 10: 17-18}{Fratres: Qui gloriátur, in Dómino gloriétur.\pscross{} Non enim qui seípsum comméndat, ille probátus est;\psstar{} sed quem Deus comméndat.}{Frères, celui qui veut être fier qu’il mette sa fierté dans le Seigneur. Celui dont on reconnaît la valeur n’est pas celui qui se recommande lui-même, c’est celui que le Seigneur recommande.}

\hymnus{MUVXLHa}
\hymnus{MUVXLHb}

\rubrica{Pour une seule vierge:}

\versiculus{Diffúsa est grátia in lábiis tuis.}{Proptérea benedíxit te Deus in ætérnum.}{La grâce est répandue sur vos lèvres.}{C'est pourquoi Dieu vous a bénie pour toujours.}

\gscore{MUVXLAB}{Ad Ben.}{Le royaume des cieux est semblable à un marchand qui cherche de belles perles ; en ayant trouvé une de grand prix, il vend tout ce qu'il a et l'achète.}

\rubrica{Pour plusieurs vierges:}

\versiculus{Adducéntur Regi Vírgines post eam.}{Próximæ ejus afferéntur tibi.}{Après elle, les vierges seront amenées au roi.}{Ses compagnes te seront amenées.}

\gscore{MUVXV1AM}{Ad Ben.}{Équipez vos lampes, vierges sages, voici l'époux vient ; sortez à sa rencontre.}

\rubrica{Oraison des Vêpres.}

\section*{À Tierce}

\capitulum{2 Co. 10: 17-18}{Fratres: Qui gloriátur, in Dómino gloriétur.\pscross{} Non enim qui seípsum comméndat, ille probátus est;\psstar{} sed quem Deus comméndat.}{Frères, celui qui veut être fier qu’il mette sa fierté dans le Seigneur. Celui dont on reconnaît la valeur n’est pas celui qui se recommande lui-même, c’est celui que le Seigneur recommande.}

\gscore{MUVXTR}{}{}

\versiculus{}{}{Dieu l’aidera par sa contemplation.
R. Dieu est en son cœur, elle ne sera pas ébranlée.}{}

\section*{À Sexte}

\capitulum{}{}{}
\gscore{}{}{}
\versiculus{}{}{Dieu l’a choisie et l’a préférée.
R. Dans son tabernacle, il la fait habiter.}{}

\section*{À None}

\capitulum{}{}{}
\gscore{}{}{}
\versiculus{}{}{}{}

\section*{À Vêpres}

\capitulum{}{}{}
\hymnus{}

\rubrica{Pour une seule vierge:}

\versiculus{}{}{}{}
\gscore{}{Ad Mag.}{}

\rubrica{Pour plusieurs vierges:}

\versiculus{Adducéntur Regi Vírgines post eam.}{Próximæ ejus afferéntur tibi.}{Après elle, les vierges seront amenées au roi.}{Ses compagnes te seront amenées.}
\gscore{MUVXV1AM}{Ad Ben.}{Équipez vos lampes, vierges sages, voici l'époux vient ; sortez à sa rencontre.}

\rubrica{Pour une vierge martyre}

\oratio{Deus, qui inter cétera poténtiæ tuæ mirácula étiam in sexu frágili victóriam martýrii contulísti:\pscross{} concéde propítius; ut, qui beátæ \rubr{N.} Vírginis et Mártyris tuæ natalícia cólimus,\psstar{} per ejus ad te exémpla gradiámur. \perdominumla}{Ô Dieu, qui, entre autres miracles de Votre puissance, avez accordé même au sexe faible la victoire du martyre, faites, dans Votre bonté, qu'honorant la naissance de Votre bienheureuse Vierge et Martyre \rubr{N.}, nous allions à Vous en imitant ses exemples. \perdominumfr}

\oratio{Indulgéntiam nobis, quǽsumus, Dómine, beáta \rubr{N.} Virgo et Martyr implóret:\pscross{} quæ tibi grata semper éxstitit, et mérito castitátis,\psstar{} et tuæ professióne virtútis. \perdominumla}{De grâce, Seigneur, que la bienheureuse Vierge et Martyre \rubr{N.} implore pour nous le pardon, elle qui Vous a toujours été agréable, et par le mérite de sa chasteté, et par la confession de Votre puissance. \perdominumfr}

\rubrica{Pour une vierge qui n'est pas martyre}

\oratio{Exáudi nos, Deus, salutáris noster:\pscross{} ut sicut de beátæ \rubr{N.} Vírginis tuæ festivitáte gaudémus;\psstar{} ita piæ devotiónis erudiámur afféctu. \perdominumla}{Exaucez-nous, ô Dieu, notre Sauveur, en sorte que, tout en nous donnant la joie, la fête de la bienheureuse \rubr{N.} votre vierge, nous instruise par le sentiment d'une pieuse dévotion.}

\rubrica{Pour plusieurs vierges martyres}

\oratio{Da nobis, quǽsumus, Dómine Deus noster, sanctárum Vírginum et Mártyrum tuárum \rubr{N.} et \rubr{N.} palmas incessábili devotióne venerári:\pscross{} ut quas digna mente non póssumus celebráre,\psstar{} humílibus saltem frequentémus obséquiis.}{Accordez-nous, Seigneur notre Dieu, de vénérer sans cesse de notre dévotion la palme de la victoire de votre saintes Vierges et Martyres \rubr{N.} et \rubr{N.}, afin que nous honorions du moins par nos humbles hommages celle que nous ne pouvons célébrer d'une âme parfaitement digne.}

%%%%%%%%%%%%%%%%%%%%%%%%%%%%%%%%%%%%%%%%%%%

\chapter*{Commun des saintes femmes}

\section*{À Laudes}

\capitulum{}{}{}
\hymnus{}
\versiculus{}{}{}{}
\gscore{}{Ad Ben.}{}
\rubrica{Oraison des Vêpres.}

\section*{À Tierce}

\capitulum{}{}{}
\gscore{}{}{}
\versiculus{}{}{}{}

\section*{À Sexte}

\capitulum{}{}{}
\gscore{}{}{}
\versiculus{}{}{}{}

\section*{À None}

\capitulum{}{}{}
\gscore{}{}{}
\versiculus{}{}{}{}

\section*{À Vêpres}

\capitulum{}{}{}
\hymnus{}
\versiculus{}{}{}{}
\gscore{}{Ad Mag.}{}
\oratio{}{}
\oratio{}{}

%%%%%%%%%%%%%%%%%%%%%%%%%%%%%%%%%%%%%%%%%%%

\chapter*{Commun de la dédicace}

\section*{À Laudes}

\capitulum{Ap. 21: 2}{Vidi civitátem sanctam, Jerúsalem novam,\pscross{} descendéntem de cælo a Deo,\psstar{} parátam sicut sponsam ornátam viro suo.}{La Ville sainte, la Jérusalem nouvelle, je l’ai vue qui descendait du ciel, d’auprès de Dieu, prête pour les noces, comme une épouse parée pour son mari.}

\hymnus{CDEDLH}

\versiculus{Hæc est domus Dómini fírmiter ædificáta.}{Bene fundáta est supra firmam petram.}{Voici la maison du Seigneur, solidement édifiée.}{Elle a été bien fondée sur le roc ferme.}

\gscore{CDEDLAB}{Ad Ben.}{Zachée, hâte-toi de descendre, car c’est dans ta maison qu’il me faut demeurer aujourd’hui. Et lui, se hâta de descendre et reçut, tout joyeux, Jésus, dans sa maison. En ce jour, à cette maison, Dieu a donné le salut, alléluia.}

\rubrica{Oraison des Vêpres.}

\section*{À Tierce}

\capitulum{Ap. 21: 2}{Vidi civitátem sanctam, Jerúsalem novam,\pscross{} descendéntem de cælo a Deo,\psstar{} parátam sicut sponsam ornátam viro suo.}{La Ville sainte, la Jérusalem nouvelle, je l’ai vue qui descendait du ciel, d’auprès de Dieu, prête pour les noces, comme une épouse parée pour son mari.}

\gscore{CDEDTR}{}{\rr A votre maison, Seigneur, convient la sainteté. \vv Tout le long des jours.}

\versiculus{Locus iste sanctus est, in quo orat sacérdos.}{Pro delíctis et peccátis pópuli.}{Ce lieu est saint, dans lequel prie le prêtre.}{Pour les délits et les péchés du peuple.}

\section*{À Sexte}

\capitulum{Ap. 21: 3}{Et audívi vocem magnam de throno dicéntem: Ecce tabernáculum Dei cum homínibus, et habitábit cum eis;\pscross{} et ipsi pópulus ejus erunt,\psstar{} et ipse Deus cum eis erit eórum Deus.}{Et j’entendis une voix forte qui venait du Trône. Elle disait : « Voici la demeure de Dieu avec les hommes ; il demeurera avec eux, et ils seront ses peuples, et lui-même, Dieu avec eux, sera leur Dieu.}

\gscore{CDEDSR}{}{\rr Ce lieu est saint, dans lequel prie le prêtre. \vv Pour les délits et les péchés du peuple.}

\versiculus{Hæc est domus Dómini fírmiter ædificáta.}{Bene fundáta est supra firmam petram.}{Voici la maison du Seigneur, solidement édifiée.}{Elle a été bien fondée sur le roc ferme.}

\section*{À None}

\capitulum{Ap. 21: 4-5}{Et abstérget Deus omnem lácrimam ab óculis eórum:\pscross{} et mors ultra non erit, neque luctus, neque clamor, neque dolor erit ultra, quia prima abiérunt.\psstar{} Et dixit qui sedébat in throno: Ecce nova fácio ómnia.}{Et Dieu essuiera toute larme de leurs yeux, et la mort ne sera plus, et il n’y aura plus ni deuil, ni cri, ni douleur : ce qui était en premier s’en est allé. Alors celui qui siégeait sur le Trône déclara : « Voici que je fais toutes choses nouvelles. »}

\gscore{CDEDNR}{}{\rr Voici la maison du Seigneur, dolidement édifiée. \vv Elle a été bien fondée sur le roc ferme.}

\versiculus{Bene fundáta est domus Dómini.}{Supra firmam petram.}{Elle est bien fondée, la maison du Seigneur.}{Sur la pierre ferme.}

\section*{À Vêpres}

\capitulum{Ap. 21: 2}{Vidi civitátem sanctam, Jerúsalem novam,\pscross{} descendéntem de cælo a Deo,\psstar{} parátam sicut sponsam ornátam viro suo.}{La Ville sainte, la Jérusalem nouvelle, je l’ai vue qui descendait du ciel, d’auprès de Dieu, prête pour les noces, comme une épouse parée pour son mari.}

\hymnus{CDEDVH}

\versiculus{Domum tuam, Dómine, decet sanctitúdo.}{In longitúdinem diérum.}{À votre maison, Seigneur, convient la sainteté.}{Tout le long des jours.}

\gscore{CDEDVAM}{Ad Mag.}{O combien il doit être révéré, ce lieu ; vraiment ce n’est pas autre chose que la maison de Dieu et la porte du ciel.}

\oratio{Deus, qui nobis per síngulos annos hujus sancti templi tui consecratiónis réparas diem, et sacris semper mystériis repræséntas incólumes:\pscross{} exáudi preces pópuli tui, et præsta; ut quisquis hoc templum benefícia petitúrus ingréditur,\psstar{} cuncta se impetrásse lætétur. \perdominumla}{Ô Dieu, qui chaque année faites revivre pour nous le jour où fut consacré ce santuaire, votre temple, et nous maintenez en vie pour participer à ces saints mystères, exaucez les prières de votre peuple, et faites que tout fidèle qui pénètre ici pour implorer vos bienfaits ait la joie de voir ses désirs comblés. \perdominumfr}

%%%%%%%%%%%%%%%%%%%%%%%%%%%%%%%%%%%%%%%%%%%

\chapter*{Commun de la Bienheureuse Vierge Marie}

\section*{À Laudes}

\capitulum{Si. 24: 9-10}{Ab inítio et ante sǽcula creáta sum,\pscross{} et usque ad futúrum sǽculum non désinam,\psstar{} et in habitatióne sancta coram ipso ministrávi.}{Dès le commencement, avant les siècles, il m’a créée, et pour les siècles je subsisterai ; dans la demeure sainte, j’ai assuré mon service en sa présence.}

\hymnus{CBMVLH}

\rubrica{À l'office de la Vierge le samedi:}
\versiculus{Benedícta tu in muliéribus.}{Et benedíctus fructus ventris tui.}{Vous êtes bénie entre les femmes.}{Et le fruit de Votre sein est béni.}
\gscore{CSMSLAB}{Ad Ben.}{Heureuse Mère de Dieu, Marie toujours Vierge, temple du Seigneur, sanctuaire du Saint-Esprit, seule plus que toute autre Vous avez plu à Jésus-Christ notre Seigneur. Priez pour le peuple, intervenez pour le clergé, intercédez pour les femmes consacrées à Dieu.}

\rubrica{Aux fêtes de la Vierge:}
\versiculus{Diffúsa est grátia in lábiis tuis.}{Proptérea benedíxit te Deus in ætérnum.}{La grâce est répandue sur vos lèvres.}{C'est pourquoi Dieu vous a bénie pour toujours.}
\gscore{A3F1VAM}{Ad Ben.}{Vous êtes bienheureuse, Marie, vous qui avez cru : car ce qui vous a été dit par le Seigneur s’accomplira en vous, alléluia.}

\oratio{Concéde nos fámulos tuos, quǽsumus, Dómine Deus, perpétua mentis et córporis sanitáte gaudére:\pscross{} et, gloriósa beátæ Maríæ semper Vírginis intercessióne,\psstar{} a præsénti liberári tristítia, et ætérna pérfrui lætítia. \perdominumla}{Seigneur notre Dieu, accordez, s’il vous plaît, à nous vos serviteurs, de jouir d’une perpétuelle santé de l’âme et du corps : et grâce à la glorieuse intercession de la bienheureuse Marie toujours Vierge, d’être délivrés des tristesses du temps présent, puis de goûter les joies éternelles. \perdominumfr}

\section*{À Prime}



\section*{À Tierce}

\capitulum{}{}{}
\gscore{}{}{}
\versiculus{}{}{}{}

\section*{À Sexte}

\capitulum{}{}{}
\gscore{}{}{}
\versiculus{}{}{}{}

\section*{À None}

\capitulum{}{}{}
\gscore{}{}{}
\versiculus{}{}{}{}

\section*{À Vêpres}

\capitulum{}{}{}
\hymnus{}
\versiculus{}{}{}{}
\gscore{}{Ad Mag.}{}
\oratio{Concéde nos fámulos tuos, quǽsumus, Dómine Deus, perpétua mentis et córporis sanitáte gaudére: et, gloriósa beátæ Maríæ semper Vírginis intercessióne, a præsénti liberári tristítia, et ætérna pérfrui lætítia.}{Seigneur notre Dieu, accordez, s’il vous plaît, à nous vos serviteurs, de jouir d’une perpétuelle santé de l’âme et du corps : et grâce à la glorieuse intercession de la bienheureuse Marie toujours Vierge, d’être délivrés des tristesses du temps présent, puis de goûter les joies éternelles.}

\end{document}